\cotefile{161-2}

\scanpage{4}

\uline{Théorème} Soient $d=(d',d'')$ un \inserted{couple de} petits ens de diagrammes, $C$ une catégorie ens. petite $\delta=(\delta',\delta'')$ un \inserted{petit} ens de cônes projectifs et inductifs dans $C$ $\Delta=\Delta_d=\Delta_{d',d''}$ $2$-Catégorie des \unsure{fatigues|catégories} ayant des lim de type $d$, avec foncteurs y commutant

\deleted{Alors il existe}

Pour tte $E\in\mathrm{Ob}\,\Delta$, soit $\underline{\mathrm{Hom}}_{\deleted{\Delta}\delta}(C,E)$ la sous-cat. \uline{pleine} de $\underline{\mathrm{Hom}}_{\deleted{\unsure{\mathrm{Ens}}}}(C,E)$ formée des foncteurs qui sont $\delta$-exacts. Si $\bar C_{\deleted{\unsure{\delta}}}$ est de même munie de $\bar\delta=(\bar\delta',\bar\delta'')$ et est $\in\Delta$, on définit $\underline{\mathrm{Hom}}_{\Delta,\delta_1}(C_1,E)$ comme intersection de $\underline{\mathrm{Hom}}_{\bar\delta}$ et $\underline{\mathrm{Hom}}_\Delta$. Ceci posé, il existe une catégorie $\bar C_{\deleted{\delta}}\in\mathrm{Ob}\,\Delta$, et un foncteur $i:C\unsure{-|\to}\bar C_{\deleted{\delta}}$, tel que

a) $i$ soit $\delta$-exact i.e. $i\in\underline{\mathrm{Hom}}_\delta(C,\bar C_{\deleted{\delta}})$

b) il soit universel en le sens suivant : Pour tt $E\in\mathrm{Ob}\,\Delta$, le foncteur $\varphi\longmapsto\varphi\circ i$ induit une équivalence

\[
\underline{\mathrm{Hom}}_{\Delta,i(\delta)}(\bar C_{\deleted{\delta}},E)\longrightarrow\underline{\mathrm{Hom}}_\delta(C,E)
\]

[ i.e. $(\bar C,i)$ est objet $2$-initial de la $2$-catégorie des $E$ \inserted{$\in\mathrm{Ob}\,\Delta$} sous $C$ tels que $C\to E$ soit $\delta$-exact ]. De plus, $\bar C_{\deleted{\delta}}$ est ens. petite.

\uline{NB} Ce Pb résoud à la fois le \uline{Pb 1} $(\unsure{\Gamma'|\delta'}=\delta''=\varnothing)$ et le \uline{Pb} 2 ($C\in\mathrm{Ob}\,\Delta$, \deleted{\illegible{1}} $\delta'=$ \deleted{diagrammes} \inserted{cônes pr. de type $d'$} exacts \deleted{\unsure{égaux à|égaux}} $\cup$ les diagrammes $(x\to y)\in\Sigma$, $\delta''=$ \deleted{diagrammes} cônes inductifs de type $d''$ exacts). On verra que le th. permet aussi de traiter le cas \deleted{\unsure{des|de}} conditions d'exactitude.

\uline{Dém} On construit \deleted{\unsure{une}} \inserted{par} récurrence transfinie \deleted{\unsure{et|une}} suite transfinie de \inserted{petites} catégories $C_i$, sous $C$, munies de $\delta_i=(\delta_i',\delta_i'')$ (petit).

\diagram{\begin{tikzcd}[ampersand replacement=\&] C\arrow[r,"\alpha_i"]\& C_i\arrow[r,"\tau_i",dash]\& C_{i+1}\end{tikzcd}}

\[
\left\{\begin{aligned}
&C_0=C,\ \delta_0=\delta=(\delta',\delta''),\\
&\deleted{(C_{i+1},\delta_{i+1})\text{ déduit de }(C_i,\delta_i)\text{ ainsi :}}\\
&(C_{i+1},\delta_{i+1})=K(C_i,\delta_i),\quad K\text{ étant explicité ci-dessous}\\
&C_{i\deleted{\unsure{+1}}}=\varinjlim_{j<i\,\mathrm{Cat}} C_i\quad\text{si }i\text{ ordinal limite},\quad
\begin{aligned}\delta_i'&=\bigcup\text{ images }\delta_j'\text{ dans } C_i\\
\delta_i''&=\bigcup\text{ im. }\delta_j''\text{ ds }C_i\end{aligned}
\end{aligned}\right.
\]

% archive sheet number: 3
